\documentclass[italian,12pt]{article}

\def\Title{Verbale interno}
\def\Subtitle{Riunione interna settimanale}
\def\Author{7Last}
\def\Date{2024-05-22}
\def\Version{v1.0}

%pacchetti extra da scaricare dblfloatfix, fancyhdr
\usepackage[left=2cm, right=2cm, bottom=3cm, top=3cm]{geometry}
\usepackage{fancyhdr}%creazione header-footer
\usepackage{graphicx} %serve per inserire immagini
\graphicspath{ {../../../logo/} }
%\usepackage{dblfloatfix} %serve per posizionare gli elementi dove si vuole
\usepackage[hidelinks]{hyperref} %serve per i link
\usepackage{tikz}
\usepackage{tgadventor} % font
\usepackage[useregional=numeric,showseconds=true,showzone=false]{datetime2}
\usepackage{caption}

\usepackage{hyperref}
\usepackage{tocloft}
\usepackage{titlesec}
\usepackage{color}
\usepackage{ulem}
\usepackage{pgfplots}
\usepackage{pgf-pie}
\usepackage[italian]{babel}
\usepackage{comment}
\usepackage{tabularx}
\usepackage{longtable}
\usepackage{float}
\usepackage{amsmath}

\linespread{1.2}
\captionsetup[table]{labelformat=empty}
\geometry{headsep=1.5cm}

\renewcommand{\contentsname}{Indice}
\renewcommand\familydefault{\sfdefault}

\begin{document}

\newgeometry{left=2cm,right=2cm,bottom=2.1cm,top=2.1cm}
\begin{titlepage}
    \vspace*{.5cm}

    \vspace{2cm}
    {
        \centering
        {\bfseries\huge \Title\par}
        \bigbreak
        {\bfseries\large \Author\par}
        \bigbreak
        {\Version\par}
        \vfill

        \begin{tikzpicture}[remember picture,overlay]

            \fill[blue!20!black] (current page.south west) -- ++(10cm,0) -- ++(-10cm,15cm) -- cycle;
            \fill[orange] (current page.south east) -- ++(-18cm,0) -- ++(21.6cm,22cm) -- cycle;

            \clip (0,-3cm) circle (2.5cm) node (current page.center) {\includegraphics[width=5cm]{logo.jpg}};
        \end{tikzpicture}

    }

    \vfill

\end{titlepage}

\restoregeometry






















\newpage

%---------------header------------------%
\pagestyle{fancy}
\fancyhead{} % pulizia degli header
\lhead{%
\begin{tikzpicture}
    \clip (0,0) circle (0.5cm);
    \node at (0,0) {\includegraphics[width=1cm]{logo}};
\end{tikzpicture}%
}
\chead{\vspace{\fill}\Title\vspace{\fill}}
\rhead{\vspace{\fill}\Version\vspace{\fill}}
%quad è una spaziatura
%---------------------------------------%

\begin{table}[!h]
	\caption{Versioni}
	\footnotesize
	\begin{center}
		\begin{tabular}{ l l l l p{6cm} }
			\hline                                                                              \\[-2ex]
			Ver. & Data       & Redattore          & Verificatore       & Descrizione           \\
			\\[-2ex] \hline \\[-1.5ex]
			1.0  & 2024-03-16 & Raul Seganfredo    & Davide Malgarise & Stesura verbale \\
			\\[-1.5ex] \hline
		\end{tabular}
	\end{center}
\end{table}

\newpage

\tableofcontents

\newpage

\section{Dettagli della riunione}


\textbf{Sede della riunione}: Google Meet\\
\textbf{Orario di inizio}: 15:30\\
\textbf{Orario di fine}: 16:30\\

\begin{flushleft}
	\begin{table}[!h]
	\begin{tabular}{ |l|l|l| } 
		\hline
		\textbf{Partecipante} & \textbf{Ruolo}       & \textbf{Presenza} \\
		\hline 
		Antonio Benetazzo     &                      & Presente          \\
		Davide Malgarise      & Verificatore         & Presente          \\
		Elena Ferro           & Amministratore       & Presente          \\
		Leonardo Baldo        &                      & Presente          \\
		Matteo Tiozzo         &                      & Presente          \\
		Raul Seganfreddo      & Redattore			 & Presente          \\
		Valerio Occhinegro    &                      & Presente          \\
		\hline
	\end{tabular}
	\end{table}
	\textbf{Partecipanti esterni}: Andrea Dorigo, Daniele Zorzi, Fabio Pallaro.\\
\end{flushleft}

\section{Ordine del giorno}
\begin{itemize}
	\item Stato di sviluppo del prodotto
	\item Discussione su eventuali modifiche
	\item Discussione su eventuale candidatura per revisione RTB
	\item Decisioni prese e conclusioni
\end{itemize}

\newpage

\section{Verbale}

\subsection{Stato di sviluppo del prodotto}
Come prima cosa, il responsabile di questo sprint, Valerio Occhinegro, ha presentato lo stato di sviluppo del prodotto,
mostrando le varie modifiche fatte. Come richiesto nel precedente incontro, sono stati illustrati i vari miglioramenti
ai grafici relativi alle dashboard, i sensori aggiunti e i grafici creati.\\
Sono inoltre state mostrate alcune modifiche fatte ai sensori delle isole ecologiche, come la modifica della
simulazione dei dati che permette loro di rimanere piene per un certo periodo di tempo e di non essere
svuotate completamente. Inoltre è stato illustrato il sistema di alerting implementato per il sensore della
temperatura, che manda delle segnalazioni in caso di temperature troppo alte.\\

\subsection{Discussione su eventuali modifiche}
Il proponente ha consigliato alcuni miglioramenti da attuare, come migliorare la leggibilità dei dati in alcuni grafici e cambiare la visualizzazione della mappa a satellitare.\\

\subsection{Discussione su eventuale candidatura per la revisione RTB}
È stata discussa con il proponente la possibilità di candidarsi per la revisione \textbf{RTB}, discutendo lo
stato di avanzamento del prodotto e le modifiche da fare per poter essere preparati. Il proponente ha
dichiara che il prodotto è pronto per la revisione al netto di alcune modifiche minori consigliate in precedenza.\\

\subsection{Decisioni prese e conclusioni}
Di comune accordo con il proponente, si è deciso di rimandare la scelta della data per il prossimo incontro,
posticipandola ad un periodo successivo alla revisione \textbf{RTB}.\\


\newpage
\begin{table}[b]
	\begin{tabular}{@{}p{5cm}p{10cm}@{}}
		Data:  & \hrulefill \\
		       &            \\
		       &            \\
		Firma del proponente: & \hrulefill \\
	\end{tabular}
\end{table}

\end{document}
